\documentclass[11pt]{article}

\usepackage[utf8]{inputenc}
\usepackage[T1]{fontenc}
\usepackage{lmodern}
\usepackage{geometry}
\usepackage{amsmath,amssymb,amsfonts,amsthm,mathtools}
\usepackage{bm}
\usepackage{enumitem}
\usepackage{microtype}
\usepackage{hyperref}
\usepackage{titlesec}
\usepackage{booktabs}
\usepackage{array}
\usepackage{setspace}
\usepackage{mathrsfs}
\usepackage{physics}

\geometry{margin=1in}
\hypersetup{colorlinks=true, linkcolor=black, urlcolor=blue, citecolor=black}
\setlist[enumerate]{topsep=0.4em,itemsep=0.35em}
\setlist[itemize]{topsep=0.3em,itemsep=0.25em}
\titleformat{\section}{\large\bfseries}{\thesection.}{0.5em}{}
\titleformat{\subsection}{\normalsize\bfseries}{\thesubsection.}{0.5em}{}
\setstretch{1.06}

\newcommand{\Aether}{\AE{}ther}
\newcommand{\Aflow}{\AE{}ther-flow}
\newcommand{\TheoryName}{\Aflow{} Interpretation of Relativity}
\newcommand{\TheoryProgram}{\Aether{} / \Aflow{} framework}
\newcommand{\Order}{\mathcal{S}}
\newcommand{\Sub}{\mathcal{A}}
\newcommand{\Ph}{\Phi}

\newtheorem{definition}{Definition}
\newtheorem{proposition}{Proposition}
\newtheorem{remark}{Remark}

\title{\textbf{The \TheoryName{}}\\[0.4em]
\large Foundations}
\author{}
\date{}

\begin{document}

\maketitle

\begin{abstract}
This paper provides the foundations statement of \TheoryName{}. The ontological core is fixed as follows: the \Aether{} is the underlying four-dimensional substrate of reality; the \Aflow{} is the intrinsic ordered motion of that substrate; observed three-dimensional space is the local experiential slice of the deeper substrate; and S-time is the experienced order of change arising from the relation between matter, light, and the \Aflow{}.

Within the broader \TheoryProgram{}, \TheoryName{} is the exact-closure theory developed in the active manuscript sequence. It retains the \Aether{} / \Aflow{} ontology while taking the effective gravitational dynamics to be exactly those of general relativity with ordinary matter coupling. The present manuscript therefore does not claim a completed microphysical derivation of Einsteinian gravity from substrate variables. Its narrower purpose is to state the original intent of the framework in disciplined language, define the ontology without loose metaphor, identify the exact relation to GR, and mark clearly the boundary between what is already adopted, what is interpreted, and what remains to be derived.

Within the flagship exact-closure package, the companion \emph{Exact-Closure Sequence Overview} is the front door and the \emph{Exact-Closure Note} is the short anchor. The present paper is the first manuscript in the five-paper modular full statement. Its task is to place ontology before derivational ambition, so that exact closure is read first and any later substrate derivation is kept explicitly secondary.
%
\end{abstract}

\section{Introduction}

The \TheoryProgram{} seeks to turn the \Aether{} / \Aflow{} ontology into a mathematically controlled theory. Within that broader program, \TheoryName{} is the exact relativistic theory developed in the active manuscript sequence. In that precise sense, it is the correct entry point for the current paper set.

\paragraph{Framework and claim boundary.}
\TheoryProgram{} denotes the broader \Aether{} / \Aflow{} framework built on the claims that the \Aether{} is the underlying four-dimensional substrate of reality, the \Aflow{} is its intrinsic ordered motion, observed three-dimensional space is the local experiential slice of that deeper substrate, S-time is the experienced order of change arising from matter, light, and the \Aflow{}, and observed expansion is the three-dimensional appearance of deeper four-dimensional ordered motion. Within that framework, \TheoryName{} denotes the exact-closure theory in which the effective gravitational dynamics are adopted to be exactly Einsteinian with universal matter coupling. In the active sequence, \emph{adoption} means use of that established relativistic dynamics without claiming substrate derivation, while \emph{derivation} is reserved for a first-principles recovery from explicit substrate variables.

The need for such an entry point is structural rather than merely editorial. The project did not begin as a deviation-first program. Its initial impulse was to propose a deeper account of what observed space, temporal ordering, and relativistic phenomenology might be \emph{about}. The present manuscript restores that order of explanation by starting from the ontology and then stating the role of \TheoryName{} inside the broader \TheoryProgram{}.

This paper adopts a strict claim boundary. It does not present a completed substrate derivation of the gravitational sector. It does not claim that ontology by itself replaces mathematical closure. Instead, it advances four narrower claims. First, the original intent of \TheoryProgram{} can be stated rigorously in terms of a four-dimensional substrate and its intrinsic ordered motion. Second, \TheoryName{} provides the exact-closure realization of that framework. Third, observed space and S-time can be described in disciplined phenomenological language without collapsing into loose metaphor. Fourth, the relation to GR can be made explicit without overstating derivational success.

The sequence role of the present manuscript should therefore be read explicitly. After the overview and the short exact-closure anchor, this foundations paper opens the modular full statement by fixing the vocabulary that the later dynamics, consistency, relativistic-recovery, and flow-geometry papers will use. In that order, exact closure is primary and derivational continuation is secondary.

\section{Framework Statement}

\begin{definition}[\TheoryProgram{}]
\TheoryProgram{} is the umbrella framework that attempts to turn the \Aether{} / \Aflow{} ontology into a full theory program with explicit kinematics, dynamics, consistency conditions, and empirical interpretation.
\end{definition}

\begin{proposition}[Active theory status]
At the present stage of the repository, \TheoryName{} is the active theory developed in full manuscript form.
\end{proposition}

This proposition fixes the conceptual order of exposition in the main manuscript sequence. The active root is organized by exact closure rather than by an unfinished low-energy alternative proposal.

\section{Original Intent}

The enduring idea behind \TheoryProgram{} is not that gravity must immediately deviate from GR. It is that the relativistic world accessible to observers may be the effective manifestation of a deeper order. Earlier formulations expressed that intuition through the image of a three-dimensional active medium. The preserved insight in that language was that time is better understood as ordered change than as an independently existing container. The weakness of that language was that it encouraged a naive picture of an ordinary three-dimensional medium moving into an external emptiness.

The present formulation reverses that picture. It treats the four-dimensional substrate as primary and the observed three-dimensional world as derivative. What embedded observers register as spatial separation, expansion, and temporal order is then interpreted as the observer-accessible appearance of deeper ordered structure. The conceptual map may be summarized as
\[
\text{\Aether{}}
=
\text{deeper four-dimensional substrate},
\]
\[
\text{\Aflow{}}
=
\text{intrinsic ordered motion of the substrate},
\]
\[
\text{observed three-dimensional space}
=
\text{local experiential slice of the deeper substrate},
\]
\[
\text{S-time}
=
\text{experienced order of change},
\]
\[
\text{observed expansion}
=
\text{three-dimensional appearance of deeper ordered motion}.
\]

The scientific importance of this reordering is methodological. Ontology is no longer allowed to masquerade as completed dynamics. Instead, ontology fixes what the framework is trying to describe, while the later mathematical manuscripts must show how that description is closed, controlled, and related to established physics.

\section{Ontology of the \Aether{} and the \Aflow{}}

\begin{definition}[\Aether{}]
The \Aether{} is the underlying four-dimensional substrate of reality posited by \TheoryProgram{}.
\end{definition}

\begin{definition}[\Aflow{}]
The \Aflow{} is the intrinsic ordered motion of the \Aether{}. It is not, without further reduction, an ordinary vector wind in observed three-dimensional space.
\end{definition}

These definitions carry both a conceptual and a formal burden. Conceptually, the \Aether{} is the deeper ontological ground rather than a substance inside a larger external box. Formally, the framework requires at least a minimal substrate arena on which ordered motion can be represented. A disciplined kinematic skeleton therefore begins by introducing a four-dimensional differentiable substrate manifold
\[
\Sub,
\qquad
\dim \Sub = 4,
\]
together with an ordered-motion map
\[
\Phi_{\lambda}:\Sub \to \Sub,
\]
where \(\lambda\) is an ordering parameter and need not be identified with physical clock time.

This notation is intentionally modest. It does not yet define the full dynamics of the substrate. It does not prove that the fundamental degrees of freedom are exhausted by a single generator field. It only records the minimum structure needed to express the claim that the substrate possesses intrinsic ordered motion. In particular, the \Aflow{} should not be reduced to a naive preferred-frame vector unless an explicit reduced model justifies that step.

\section{Why the \Aflow{} Is Not a Naive Vector Wind}

The flow language of the framework must be handled under strict control. If the \Aflow{} were treated as nothing more than an ordinary vector wind in observed space, the framework would risk importing precisely the wrong kind of low-energy structure: spurious preferred-frame signals, the wrong symmetry interpretation, and an overly literal picture of the ontology.

For that reason, the present manuscript uses flow language in a restricted sense. The \Aflow{} names the intrinsic ordered motion of the deeper substrate, not a completed low-energy observable field content. Vectorial or congruence-like representations may still be useful as reduced kinematic proxies, but they should be read as derived descriptions, not as the full ontological content of the theory.

This discipline is especially important in \TheoryName{}. The exact-closure theory retains the ontological role of the \Aflow{} while denying that it has already appeared as an extra established weak-field or radiative sector beyond GR.

\section{Observed Space as Local Experiential Slice}

Observed space is not identified with the whole ontology of the framework. Instead, it is the local observer-accessible slice of the deeper substrate.

To state that relation formally, let
\[
\iota_s:\Sigma_s \hookrightarrow \Sub
\]
denote an observer-accessible three-dimensional slice labeled by an order parameter \(s\). If the substrate carries a geometric or proto-geometric structure \(G_{AB}\), then the induced spatial geometry on the slice is represented schematically by
\[
\gamma_{ij}(s,x) = \iota_s^{*} G_{AB}.
\]
The point of this expression is not to claim a finished emergent-geometry theorem. It is to fix the direction of explanation. Observed three-dimensional spatial relations are treated as induced or experienced structure, not as the deepest available level of description.

This is why the phrase \emph{experienced space} can be tolerated only as secondary explanatory language. The stronger formal phrase is \emph{observed three-dimensional space as a local experiential slice of the deeper substrate}. That wording makes clear that the observed world is real as phenomenology and measurement, but is not taken to exhaust the ontology.

\section{S-Time as Experienced Order of Change}

The concept of S-time is meant to preserve the earliest intuition of the program while removing its more unstable imagery. S-time is not a place-like extra corridor through which the world moves. It is the experienced order of change arising from the relation between matter, radiation, observers, and the substrate organization.

A minimal formal statement is to treat S-time as an order functional,
\[
S = \Order[\Gamma,\Psi,\Sub,\Phi_{\lambda}],
\]
defined on an observer-history \(\Gamma\), a physical state \(\Psi\), the substrate \(\Sub\), and the ordered-motion map \(\Phi_{\lambda}\). The required directional property is monotonicity along admissible histories,
\[
\frac{dS}{d\lambda} > 0.
\]
This formulation does not yet identify the final microphysical definition of S-time. It does, however, state the exact conceptual role of the term. Time in the present framework is the experienced order of physical change, not an independently existing container that precedes all dynamics.

The phrase \emph{experienced time} may therefore be used only as shorthand for S-time. Whenever such prose appears, it should be reduced immediately back to the formal statement above.

\section{Observed Expansion, Light, and Relativistic Phenomenology}

The framework does not deny the observational reality of expansion, clock behavior, redshift, or relativistic duration effects. It relocates their explanatory status. What is primary in the ontology is the deeper ordered substrate; what is primary in observation is the effective relativistic structure accessible to embedded observers.

The conceptual relation may be summarized schematically as
\[
\begin{aligned}
\text{primary ontology: } (\Sub,\Phi_{\lambda})
\qquad &\Longrightarrow \\
&\text{observed appearance: evolving slice geometry,} \\
&\text{light propagation, and temporal ordering}.
\end{aligned}
\]

On this view, observed expansion is not rejected; it is interpreted as the three-dimensional appearance of deeper four-dimensional ordered motion. Likewise, the empirical role of light remains non-negotiable. Any successful theory of \TheoryProgram{} must recover an effective relativistic structure in which null propagation, proper-time measurement, time dilation, and length-contraction phenomenology are all represented correctly in the regimes where they are tested.

The present foundations paper does not claim that those phenomena have already been derived from substrate microphysics. Its narrower claim is that the ontology is compatible with their exact recovery and that, in \TheoryName{}, that recovery is secured by exact closure to GR.

\section{\TheoryName{} as the Exact-Closure Realization of \TheoryProgram{}}

\begin{definition}[\TheoryName{}]
\TheoryName{} is the exact-closure, GR-consistent theory of \TheoryProgram{} in which the \Aether{} / \Aflow{} ontology is retained while the effective gravitational dynamics are taken to be exactly those of general relativity with ordinary matter coupling.
\end{definition}

This theory is exact closure rather than deformation. Its effective action is
\begin{equation}
S_{\mathrm{eff}}
=
\frac{c^3}{16\pi G}\int d^4x\,\sqrt{-g}\,R
+
S_{\mathrm{matter}}[g,\psi],
\label{eq:SA}
\end{equation}
with field equations
\begin{equation}
G_{\mu\nu} = \frac{8\pi G}{c^4} T_{\mu\nu}.
\label{eq:einstein}
\end{equation}
Null propagation and proper-time measurement are therefore governed by the standard relativistic relations
\begin{equation}
0 = g_{\mu\nu}\,dx^{\mu}dx^{\nu},
\label{eq:null}
\end{equation}
\begin{equation}
d\tau^2 = -\frac{1}{c^2} g_{\mu\nu}\,dx^{\mu}dx^{\nu}.
\label{eq:propertime}
\end{equation}

\begin{proposition}
The predictive content of \TheoryName{} is exactly that of GR. Accordingly, \TheoryName{} carries no independent low-energy non-GR observable signature.
\end{proposition}

\begin{remark}
This does not make \TheoryName{} scientifically empty. It makes its status precise. The theory is an exact ontological completion of GR within \TheoryProgram{}, not a distinct low-energy alternative to it.
\end{remark}

\section{\TheoryName{} as Benchmark and Reversion Theory}

\begin{proposition}[Benchmark role]
Any future extension of \TheoryProgram{} must recover \TheoryName{} in the appropriate limit. \TheoryName{} therefore functions as the internal exact benchmark of the full framework.
\end{proposition}

\begin{proposition}[Automatic reversion]
If a proposed extension of the framework fails an essential consistency or viability gate, the scientific status of the framework contracts back to \TheoryName{} rather than collapsing into incoherence.
\end{proposition}

This benchmark structure is one of the central reasons to make \TheoryName{} the active manuscript line now. A disciplined theory program cannot allow an unfinished deviation sector to remain the narrative center once the exact benchmark theory is already available and the currently tested deviation theory has not earned promotion.

The correct short status statement is therefore the following: \TheoryName{} is the exact-closure theory developed in the present paper and is the active reference theory of the repository.

\section{Relation to General Relativity}

The relation between \TheoryName{} and GR is exact at the operational level and interpretive at the ontological level. Operationally, GR remains the governing relativistic theory of gravitation in this theory. Ontologically, \TheoryProgram{} adds the proposal that the relativistic metric description may be the effective manifestation of a deeper ordered substrate.

Two clarifications are essential. First, \TheoryName{} does not dispute any tested relativistic phenomenon. Clock behavior, null propagation, redshift, time dilation, and weak-field structure are all inherited from GR exactly through \eqref{eq:SA}--\eqref{eq:propertime}. Second, the preservation of the \Aether{} / \Aflow{} ontology does not license the introduction of an extra measured preferred-frame field in the effective theory. In \TheoryName{}, the ontology acts as deeper interpretive structure, not as an additional established low-energy sector.

Local special-relativistic kinematics are also retained through this exact closure. Since the effective metric sector is Einsteinian, the local inertial description continues to reproduce the standard special-relativistic content seen by observers. The full relation to SR should still be treated more explicitly in a later relativistic-recovery manuscript, but the foundations-level point is already clear: the deeper ontology is not offered as a contradiction of relativity, but as a possible deeper account of why an effective relativistic structure governs observation.

\section{Claim Boundary and Program Status}

The present paper separates five levels of statement that should not be blurred together.

\subsection{Ontology}

The framework claims that the \Aether{} and the \Aflow{} provide the correct conceptual ground of the theory, that observed three-dimensional space is a local experiential slice of that deeper substrate, and that S-time is experienced order of change.

\subsection{Exact closure}

The framework claims that \TheoryName{} provides an exact relativistic closure in which the operational content of GR is recovered without modification.

\subsection{Adopted dynamics}

The framework adopts the Einstein equations as the effective dynamics of gravity in \TheoryName{}. This is an adoption claim and should not be misdescribed as a completed substrate derivation.

\subsection{Open derivational work}

What remains unfinished is the recovery of the effective relativistic sector from explicit substrate variables, a closed substrate action, controlled matter coupling, mode analysis, and consistency proofs. Those tasks belong to the later dynamics and consistency manuscripts.

\subsection{Established in the present paper}

The following points are claimed as established at the present foundations level:
\begin{itemize}
    \item the basic ontology of \Aether{}, \Aflow{}, observed three-dimensional space, and S-time;
    \item the theory-level definition of \TheoryName{} as the exact-closure realization of \TheoryProgram{};
    \item the exact operational relation of \TheoryName{} to GR;
    \item the benchmark and reversion role of \TheoryName{} inside the total framework.
\end{itemize}

\subsection{Not yet established}

The following stronger statements are \emph{not} claimed here:
\begin{itemize}
    \item a completed substrate derivation of Einsteinian dynamics;
    \item a full microphysical definition of the substrate degrees of freedom;
    \item a completed consistency analysis of a specific substrate action;
    \item a completed alternative low-energy sector beyond the present exact-closure line.
\end{itemize}

\section{Conclusion}

This foundations manuscript now provides the opening module of the fuller public-facing statement of \TheoryProgram{}. Its central claims are limited but definite. The \Aether{} is the deeper four-dimensional substrate of reality. The \Aflow{} is its intrinsic ordered motion. Observed three-dimensional space is the local experiential slice of that deeper substrate. S-time is the experienced order of change arising from the relation between matter, light, and the \Aflow{}. Within that ontology, \TheoryName{} is the exact-closure theory whose effective gravitational dynamics are exactly those of GR.

The result is a stable starting point rather than a completed final theory. It tells the reader what the \TheoryProgram{} is about, what \TheoryName{} is, why the \Aflow{} should not be read as a naive vector wind, how the ontology relates to observed experience, why \TheoryName{} is the active reference theory, and what the framework presently adopts rather than derives. That is the proper foundations layer for the modular full statement that continues through dynamics, consistency, relativistic recovery, and flow geometry. Read in that flagship sequence, exact closure is fixed first, while any later derivational continuation remains a downstream and explicitly secondary burden.

\end{document}
